% =============================================================================
% LC-020 - Valid Parentheses
% =============================================================================
% STATUS: draft
% PROMOTE: exemplars
% TAGS: stacks, parsing
% DIFFICULTY: easy
% SOURCE: LeetCode 020
% =============================================================================

\ProblemTitle{Valid Parentheses}
\ProblemMeta{Easy}{Stacks}{LC-020}
\label{prob:stacks-queues-valid-parentheses}

\ProblemSummary
Given a string containing only the characters \texttt{()[]\{\}}, determine whether it is valid.
A string is valid if every opening bracket is closed by the same type of bracket, and brackets
are closed in the correct order.

\KeyIdea
Use a stack to represent the sequence of currently-unmatched opening brackets.
When an opening bracket is seen, push it. When a closing bracket is seen, it must match the
most recent unmatched opening bracket (the stack top). Any mismatch, or attempting to close
when the stack is empty, makes the string invalid. The string is valid iff the stack is empty
at the end.

% -----------------------------------------------------------------------------
% Optional: Author Thinking (excluded from exemplar builds)
% -----------------------------------------------------------------------------
\begin{Thinking}
My first attempt jumped into coding the stack loop before writing the invariant. I was
conceptually correct (LIFO matching), but I missed two edge cases that matter in C++:
(1) calling \texttt{top()} on an empty stack (crash/UB), and (2) ensuring we \texttt{pop()}
after a successful match (easy to accidentally skip with \texttt{continue}).
Writing the invariant first made the implementation obvious:
``stack = exactly the unmatched openers so far''.
\end{Thinking}

% -----------------------------------------------------------------------------
% Algorithm
% -----------------------------------------------------------------------------
\Algorithm
\begin{CodeBlock}[style=pseudo,caption={Pseudocode}]{12}
stack = empty
fo c in s:
    if c is '(' or '[' or '{':
        push c onto stack
    else:
        if stack is empty: return false
        t = top of stack
        if (c,t) is not one of (')','('), (']','['), ('}','{'):
            return false
        pop stack
return (stack is empty)
\end{CodeBlock}

% -----------------------------------------------------------------------------
% Correctness
% -----------------------------------------------------------------------------
\Correctness
Maintain the invariant: after processing any prefix of the string, the stack contains exactly
the opening brackets from that prefix that have not yet been matched, in the order they were
encountered.
When an opening bracket is read, pushing it preserves the invariant by adding a new unmatched
opener. When a closing bracket is read, validity requires that there exists an unmatched opener
and that the most recent unmatched opener matches the closing type; this is exactly the stack
top. If the stack is empty or the types mismatch, the string cannot be valid, and the algorithm
returns false. Otherwise, popping removes the matched opener and preserves the invariant.
After the full scan, the string is valid if no unmatched openers remain, i.e., the stack is empty.
Therefore, the algorithm returns true exactly for valid strings.

% -----------------------------------------------------------------------------
% Complexity
% -----------------------------------------------------------------------------
\Complexity
Let $n$ be the length of the string. Each character is pushed and popped at most once, so the
runtime is $O(n)$. The stack stores at most $n$ opening brackets, so the extra space is $O(n)$.

% -----------------------------------------------------------------------------
% Implementation
% -----------------------------------------------------------------------------
\bigskip
\noindent\textbf{C++ Implementation}

\lstinputlisting[
  style=cpp,
  caption={C++ Solution: Valid Parentheses}
]{../problems/05-stacks-queues/easy/lc-020-valid-parenthesessolution.cpp}

\bigskip
\noindent\textbf{Python Implementation}

\lstinputlisting[
  language=Python,
  backgroundcolor=\color{codebg},
  basicstyle=\ttfamily\small,
  frame=single,
  breaklines=true,
  caption={Python Solution: Valid Parentheses}:\While{}
    \State 
  \EndWhile
]{../problems/05-stacks-queues/easy/lc-020-valid-parenthesessolution.py}

\ProblemDivider
